\section{Compositionality of Secure Compilers}

\begin{definition}[Universal Image]\label{def:sigma}
  \[
  \sigma_{\sim}(\trg{\pi}) = \left\{ \src{\trace} | \forall\trg{\trace},\src{\trace}\sim\trg{\trace}\implies\trg{\trace}\in\trg{\pi}\right\}
\]
\end{definition}
\begin{definition}[Existential Image]\label{def:tau}
  \[
  \tau_{\sim}(\src{\pi}) = \left\{ \trg{\trace} | \exists\src{\trace},\src{\trace}\sim\trg{\trace}\wedge\src{\trace}\in\src{\pi}\right\}
\]
\end{definition}

\begin{definition}[Robust Trace-Hyperproperty Preservation with Universal Image]\label{def:rtp:sigma}
  For a given class $\trg{\cC}$, a compiler from languages $\S$ to $\T$ robustly preserves $\trg{\cC}$ iff
  $$
  \forall\trg{\Pi}\in\trg{\cC},\forall\src{p}\in\src{\partials},\rsat{\src{p}}{\sigma_{\sim}(\trg{\Pi})}\implies\rsat{\stcomp{p}}{\trg{\Pi}}
  $$
  We write $\rtpsigma{\stcomp{\bullet}}{\trg{\cC}}$.
\end{definition}


\begin{definition}[Robust Trace-Hyperproperty Preservation with Existential Image]\label{def:rtp:tau}
  For a given class $\cC$, a compiler from languages $\S$ to $\T$ robustly preserves $\src{\cC}$ iff
  $$
  \forall\src{\Pi}\in\src{\cC},\forall\src{p}\in\src{\partials},\rsat{\src{p}}{\src{\Pi}}\implies\rsat{\stcomp{p}}{\tau_{\sim}(\src{\Pi})}
  $$
  We write $\rtptau{\stcomp{\bullet}}{\src{\cC}}$.
\end{definition}

\begin{definition}[Robust Trace-Hyperproperty Preservation]\label{def:rtp}
  For a given class $\cC$, a compiler from languages $\S$ to $\T$ robustly preserves $\cC$ iff
  $$
  \forall\Pi\in\cC,\forall\src{p}\in\src{\partials},\rsat{\src{p}}{\Pi}\implies\rsat{\stcomp{p}}{\Pi}
  $$
  We write $\rtp{\stcomp{\bullet}}{\cC}$.
\end{definition}

\begin{definition}[Sequential Composition of Compilers]
  Given two compilers $\sicomp{\bullet}$ and $\itcomp{\bullet}$, their sequential composition is $\sitcomp{\bullet}=\itcompN{\sicomp{\bullet}}$.
\end{definition}
% we can propagate through the assumptions from src langs to intermediate langs
%

\begin{lemma}[Weakening RTP]\label{lem:weaken}
  Given classes $\cC_{1}, \cC_{2}$ such that
  \begin{assumptions}
    \item $\cC_{1}\subseteq\cC_{2}$
    \item $\rtp{\stcomp{\bullet}}{\cC_{2}}$
  \end{assumptions}
  We show
  \begin{goals}
    \item $\rtp{\stcomp{\bullet}}{\cC_{1}}$
  \end{goals}
\end{lemma}
\begin{proof}
  Using \Cref{def:rtp} on the goal, let $\Pi\in\cC_{1}$ and $\src{p}\in\src{\partials}$ such that $\rsat{\src{p}}{\Pi}$, so what's left to prove is $\rsat{\stcomp{p}}{\Pi}$.
  Since $\cC_{1}\subseteq\cC_{2}$ and $\Pi\in\cC_{1}$, we know that $\Pi\in\cC_{2}$.
  Thus, we can apply the assumption $\rtp{\stcomp{\bullet}}{\cC_{2}}$ to our goal, leaving us with $\rsat{\src{p}}{\Pi}$ to show, which was an assumption we made.
\end{proof}

\begin{definition}[Well-formedness of $\sim$ for a Class $\src{\cC}$]\label{def:wfc:tau:tracerel}
  % $$

  \[
    \wfctau{\sim}{\src{\cC}} := \forall \src{\pi}\in\src{\cC}, \tau_{\sim}(\src{\pi})\in\tau_{\sim}(\src{\cC})
  \]
  % $$
\end{definition}
\begin{definition}[Well-formedness of $\sim$ for a Class $\trg{\cC}$]\label{def:wfc:sig:tracerel}
  % $$

  \[
    \wfcsig{\sim}{\trg{\cC}} := \forall \trg{\pi}\in\trg{\cC}, \sigma_{\sim}(\trg{\pi})\in\sigma_{\sim}(\trg{\cC})
  \]
  % $$
\end{definition}

\begin{definition}[Safety Properties]
  The class of safety properties contains the lifting of all properties that can be refuted with a finite trace prefix:
  $$
  \cSafety = \left\{\lift{\prop}\ |\ \forall \trace\in\traces, t\not\in\lift{\prop} \text{ iff } \exists m\ge\trace,\forall \trace'\in\traces,m\le\trace'\implies\trace'\not\in\lift{\prop}\right\}
  $$
\end{definition}

\begin{definition}[Hypersafety Properties]\label{def:hsafety}
  The class of hypersafety properties contains all hyperpropert that can be refuted with an observation:
  $$
  \cHSafety = \left\{\Pi\ |\ \forall b\in 2^{\traces},b\not\in\Pi\text{ iff  }\exists o\ge b,\forall b'\in 2^{\traces},o\le b'\implies b'\not\in\Pi\right\}
  $$
\end{definition}

\begin{definition}[Subset Closed Hyperproperties]
  The class of hyperproperties that are closed with respect to the subset relation is
  $$
  \cSS = \left\{H\ |\ \forall X\in H, \forall Y\subseteq X, Y \in H\right\}
  $$
\end{definition}

\begin{lemma}[Hypersafety is entailed in SSC]
  $\cHSafety\subseteq\cSS$.
\end{lemma}
\begin{proof}
  \cite{clarkson08}
\end{proof}

\begin{definition}[K-Hypersafety]
  Exactly the same as \Cref{def:hsafety}, but the observations $o$ are restricted to cardinality $k$.
  2-Hypersafety is simply $k=2$. \Cref{def:ni} gives an example instance of a classic 2-hypersafe property.
\end{definition}

\begin{definition}[Non-Interference ($\Ni$)]\label{def:ni}
  We define the class containing the non-interference hyperproperty as:
  $$
  \Ni = \left\{ H | \forall \trace_{1},\trace_{2}\in H. \loweq{\trace_{1}}{\trace_{2}}\implies\trace_{1}=\trace_{2} \right\}
  $$
\end{definition}
Note that $=$ may not be strict equality, but some suitable trace equivalence that checks both public and private actions, instead of just public.

